\begin{center}
    \textbf{Abstract}
\end{center}

\noindent The use of prestressed concrete structures grew and gained prominence after World War II due to the development of high-strength steel, given their constructive versatility for structures requiring larger free spans. Among the calculations involved in designing projects using this construction method are prestress losses, which pose a real challenge due to their scope and complexity. Although the precise calculation of these losses is crucial for structural optimization, it is often neglected, resulting in over-dimensioning and increased costs. Existing computational tools, often opaque and lacking transparency, do not adequately address this aspect, leading to the use of simplified constants. Given this scenario, the present study proposes the development of an accessible, free, and open-source computational tool written in the Python programming language to facilitate the transparent and modified calculation of prestress losses, promoting clarity, efficiency, and economy in the design of prestressed concrete structures.

\vspace{0.2cm}

\noindent\textbf{Keywords}: prestressed concrete; prestress loss; computational tool.